\documentclass[11pt]{amsart}
\usepackage{amsmath,amssymb,amsthm}
\usepackage{graphicx}
\usepackage{hyperref}

\newtheorem{theorem}{Theorem}
\newtheorem{corollary}[theorem]{Corollary}

\newcommand{\Cp}{C_p}
\newcommand{\spn}{\operatorname{span}}

\title{\(C_p\)-Semi-embeddings Always Embed}
\author{Agent lane 05}
\date{June 15, 2026}

\begin{document}
\maketitle

\section*{Source And Target}

The source paper is Kakol--Kurka--Sliwa, \emph{Lotz-Peck-Porta and
Rosenthal's theorems for spaces \(C_p(X)\)}, arXiv:2509.08634 \cite{KKS}.
On page 4, Problem 10 asks:

\begin{quote}
Let \(X\) be an infinite compact space. Is it true that each semi-embedding
\(T:\Cp(X)\to \Cp(Y)\) is an embedding for any compact space \(Y\) if and
only if \(X\) is scattered?
\end{quote}

\begin{center}
\includegraphics[width=0.96\linewidth]{figures/open_problem_crop.png}
\end{center}

Here \(\Cp(X)\) denotes \(C(X)\) with the topology of pointwise convergence.
Following the source paper, a continuous linear map \(T:E\to F\) between
locally convex spaces is \emph{polar} if there is a base \(\mathcal U\) of
closed zero-neighborhoods in \(E\) such that \(T(U)\) is closed in \(F\) for
every \(U\in\mathcal U\).  A \emph{semi-embedding} is an injective polar map.

\section*{Result}

Problem 10 has a negative answer.  The proposed ``if and only if \(X\) is
scattered'' characterization is false, because the semi-embedding property
does not detect scatteredness at all: between \(C_p\)-spaces, every
semi-embedding is automatically a topological embedding.

\section*{Proof Intuition}

A continuous linear map \(T:\Cp(X)\to\Cp(Y)\) is controlled point by point on
the codomain.  For each \(y\in Y\), the functional
\(\delta_y\circ T\) on \(\Cp(X)\) depends on only finitely many points of
\(X\).  These finite-support functionals generate the topology inherited by
the range \(T(\Cp(X))\).  If \(T\) were not an embedding, some point evaluation
\(\delta_x\) on \(X\) would not be continuous for this inherited topology.
The polar condition rules this out.  Indeed, take a polar neighborhood whose
preimage controls \(f(x)\).  It contains a finite-coordinate kernel
\(M=\{f:f|_F=0\}\), and its closedness in the inherited topology forces the
closure of \(M\) to remain inside the strip \(|f(x)|\le 1\).  Since that
closure is a linear subspace, \(\delta_x\) annihilates it.  The annihilator of
that closure is exactly the part of the inherited dual supported on the finite
set \(F\), so \(\delta_x\) must already belong to the inherited dual.

\begin{theorem}
Let \(X\) and \(Y\) be Tychonoff spaces.  Every semi-embedding
\[
  T:\Cp(X)\longrightarrow \Cp(Y)
\]
is a topological embedding.
\end{theorem}

\begin{proof}
Let \(E=C(X)\) as a vector space, and let
\(\tau_p=\sigma(E,\{\delta_x:x\in X\})\) be the topology of \(\Cp(X)\).
For each \(y\in Y\), put
\[
  h_y=\delta_y\circ T\in E'.
\]
Since every continuous linear functional on \(\Cp(X)\) is a finite linear
combination of point evaluations, each \(h_y\) belongs to
\(\spn\{\delta_x:x\in X\}\).  Let
\[
  H=\spn\{h_y:y\in Y\}\subset E'
\]
and let \(\rho=\sigma(E,H)\).  The topology \(\rho\) is precisely the topology
pulled back from the subspace \(T(E)\subset \Cp(Y)\).  Since \(T\) is
injective, \(T:(E,\rho)\to T(E)\) is a linear homeomorphism onto its range.
Also \(\rho\le \tau_p\), because \(T:\Cp(X)\to\Cp(Y)\) is continuous.

We claim that \(\delta_x\in H\) for every \(x\in X\).  Fix \(x_0\in X\).
Because \(T\) is polar, there is a base \(\mathcal U\) of closed
\(\tau_p\)-neighborhoods such that \(T(U)\) is closed in \(\Cp(Y)\) for every
\(U\in\mathcal U\).  Choose \(U\in\mathcal U\) with
\[
  U\subseteq \{f\in E: |f(x_0)|\le 1\}.
\]
The set \(T(U)\) is closed in \(\Cp(Y)\), hence closed in the range \(T(E)\).
Transporting this statement through the homeomorphism
\((E,\rho)\cong T(E)\), we get that \(U\) is \(\rho\)-closed.

Since \(U\) is a \(\tau_p\)-neighborhood of zero, there are a finite set
\(F\subset X\) and \(\varepsilon>0\) such that
\[
  V=\{f\in E: |f(t)|<\varepsilon\text{ for all }t\in F\}\subseteq U.
\]
Let
\[
  M=\{f\in E: f(t)=0\text{ for all }t\in F\}.
\]
Then \(M\subset V\subset U\).  Since \(U\) is \(\rho\)-closed, the
\(\rho\)-closure \(\overline M^\rho\) is contained in \(U\).  The set
\(\overline M^\rho\) is a linear subspace.  Therefore, if
\(g\in\overline M^\rho\), then \(ag\in\overline M^\rho\subset U\) for every
scalar \(a\).  As \(U\subseteq\{|f(x_0)|\le 1\}\), this implies
\(|a g(x_0)|\le 1\) for every scalar \(a\), hence \(g(x_0)=0\).  Thus
\[
  \delta_{x_0}\in(\overline M^\rho)^\perp .
\]

It remains to identify this annihilator.  In the weak topology \(\rho\), the
closure of \(M\) is
\[
  \overline M^\rho
  =\bigcap_{\lambda\in H\cap M^\perp}\ker \lambda .
\]
The annihilator of this intersection is the pointwise closure of
\(H\cap M^\perp\) in the algebraic dual.  But
\[
  M^\perp=\spn\{\delta_t:t\in F\}
\]
is finite-dimensional, so the subspace \(H\cap M^\perp\) is pointwise closed.
Hence
\[
  (\overline M^\rho)^\perp=H\cap M^\perp
\]
as subspaces of the algebraic dual.  Consequently \(\delta_{x_0}\in H\).
Since \(x_0\) was arbitrary, every point evaluation \(\delta_x\) is
\(\rho\)-continuous.

The topology \(\tau_p\) is generated by the point evaluations \(\delta_x\).
Thus \(\tau_p\le \rho\).  We already had \(\rho\le\tau_p\), so
\(\rho=\tau_p\).  Therefore the range topology inherited from \(\Cp(Y)\) is
exactly the original topology of \(\Cp(X)\), and \(T\) is a topological
embedding.
\end{proof}

\begin{corollary}
The characterization proposed in Problem 10 of \cite{KKS} is false.
\end{corollary}

\begin{proof}
Take \(X=[0,1]\).  This is an infinite compact non-scattered space.  By the
theorem, for every compact space \(Y\), each semi-embedding
\(T:\Cp([0,1])\to\Cp(Y)\) is an embedding.  Thus the property appearing in
Problem 10 holds although \(X\) is not scattered.
\end{proof}

\section*{Verification Notes}

No computational check is involved.  The proof uses only:
\begin{itemize}
\item the standard description of \(\Cp(X)'\) as finite linear combinations
of point evaluations;
\item the fact that \(T(U)\) closed in \(\Cp(Y)\) makes \(U\) closed for the
range-induced topology \(\rho=\sigma(E,H)\);
\item the closure/annihilator formula for a subspace in a weak topology.
\end{itemize}

The argument also explains why the dense-restriction map mentioned in the
source paper is not a semi-embedding: the induced weaker topology does not
make the missing point evaluations continuous, so it cannot have the required
closed-neighborhood base.

\section*{Novelty And Search Bounds}

Before promotion, the run's registry, solution, attempt, and proof-gap indexes
were searched for arXiv:2509.08634, the title keywords, and the core phrases
``semi-embedding'', ``\(C_p(X)\)'', ``\(C_p(Y)\)'', and ``polar map''.  No
existing packet for this result was found.

External web searches on 2026-06-15 used exact and near-exact phrases from
Problem 10, including
``Let X be an infinite compact space semi-embedding \(C_p(X)\) \(C_p(Y)\)'',
``semi-embedding \(C_p(X)\) \(C_p(Y)\) compact scattered'', and
``polar map \(C_p(X)\) quasibarrelled semi-embedding''.  No later paper
explicitly answering Problem 10 was found.  The result should nevertheless be
reviewed with medium novelty confidence, because the proof is short and may
exist as folklore in \(C_p\)-theory.

\section*{Human Review Recommendation}

Review as a candidate full negative answer to Problem 10.  The key checks are:
first, that the source definition of semi-embedding allows the cofinal closed
neighborhood \(U\) used above; second, that the induced topology on \(T(E)\)
is exactly \(\sigma(E,H)\); and third, that the double-annihilator step is
valid for the weak closure of \(M\), using the finite-dimensionality of
\(M^\perp\).

\begin{thebibliography}{9}
\bibitem{KKS}
J. Kakol, O. Kurka, and W. Sliwa,
\emph{Lotz-Peck-Porta and Rosenthal's theorems for spaces \(C_p(X)\)},
arXiv:2509.08634, 2025.

\bibitem{Ark}
A. V. Arkhangel'skii,
\emph{Topological Function Spaces},
Kluwer Academic Publishers, 1992.
\end{thebibliography}

\end{document}
